% !TeX root = ../../Book4.tex
% 纲要标题：### 6.1 复形
% 纲要标签：b4:sec:0501
% 纲要位置：计划纲要.md，第 3218 行。

\section{复形}
\label{b4:sec:0501}

相邻映射的复合为零，只保证像包含在核中。同调把两者的差额组织成逐次数的不变量，链与上链的编号则须先固定转换规则。

% 纲要内容（仅作写作计划，尚非教材正文）：
%
% * 链复形与上链复形
% * 圈（cycle）、边界（boundary）及同调
% * 指标转换与次数平移约定
%

\subsection{链复形与上链复形}
\label{b4:subsec:0501:01}

一个序列若处处正合，核与像就完全一致；但许多计算只能保证相邻两步的复合为零。先保留这一较弱的条件，才有可能记录正合性在哪个次数失败，以及核中还剩下什么。

\begin{definition}\label{b4:def:cochain-complex}
设 \(\mathcal A\) 是交换范畴。给定对象 \(C^n\) 及态射 \(\mathrm{d}_C^n:C^n\to C^{n+1}\)（\(n\in\mathbb Z\)），若每个 \(n\) 都满足 \(\mathrm{d}_C^{n+1}\mathrm{d}_C^n=0\)，则称 \(C=(C^n,\mathrm{d}_C^n)\) 为 \(\mathcal A\) 中的\term[shanglian-fuxing]{上链复形}[cochain complex]，\(\mathrm{d}_C^n\) 称为微分。类似地，对象 \(C_n\) 及态射 \(\partial_n:C_n\to C_{n-1}\) 满足 \(\partial_{n-1}\partial_n=0\) 时，组成\term[lian-fuxing]{链复形}[chain complex]。
\end{definition}
\symidx{\(\mathrm{d}_C^n,\ \partial_n\)：上链复形与链复形的微分}

复形条件在每个次数只保证
\[
\mathrm{d}_C^n\mathrm{d}_C^{n-1}=0
\quad\Longleftrightarrow\quad
\operatorname{im}\mathrm{d}_C^{n-1}\subseteq\ker\mathrm{d}_C^n;
\]
正合还要求这里的像与核是同一个子对象。同调将记录这个包含还差多少。链与上链的差别在指标方向：链微分降一次，上链微分升一次。它们的核、像和正合性问题相同。为统一后面的 Hom 复形、移位和映射锥，本编主要用上链指标；讨论消解或 Tor 时也会使用链指标，并明确作转换。

\begin{example}\label{b4:exm:two-term-complex}
在左 \(R\)-模范畴中，任一同态 \(u:M\to N\) 都给出仅在次数 \(0,1\) 非零的复形 \(M\xrightarrow{u}N\)。此时 \(\mathrm{d}^0=u\)，下一微分是 \(\mathrm{d}^1=0_{N,0}:N\to0\)，所以所需条件为
\[
\mathrm{d}^1\mathrm{d}^0=0_{N,0}\circ u=0.
\]
其余相邻复合也均为零。不要求 \(u^2=0\)，因为 \(u\) 的目标是 \(N\)，而源是 \(M\)，一般根本不能把 \(u\) 与自身复合。相反，若每个次数均放同一模 \(M\)，并在每一步都用同一自同态 \(a\)，便确实需要 \(a^2=0\)。例如 \(R=k[\varepsilon]/(\varepsilon^2)\) 上反复乘 \(\varepsilon\) 给出复形。
\end{example}

\begin{definition}\label{b4:def:bounded-complex}
设 \(C\) 是交换范畴中的上链复形。若存在整数 \(a\)，使 \(n<a\) 时 \(C^n=0\)，则称 \(C\) 下有界；若存在 \(b\)，使 \(n>b\) 时 \(C^n=0\)，则称其上有界；同时满足二者时称其\term[youjie-fuxing]{有界复形}[bounded complex]。对象 \(M\) 放在次数零、其余项为零所得复形记为 \(M[0]\)。
\end{definition}
\symidx{\(M[0]\)：集中在零次的复形}

这里的上下指整数次数的大小，与箭头在版面上朝哪个方向无关。

\subsection{圈、边界与同调}
\label{b4:subsec:0501:02}

\begin{definition}\label{b4:def:cohomology}
设 \(C\) 是交换范畴 \(\mathcal A\) 中的上链复形。记
\[
Z^n(C)\coloneqq\ker \mathrm{d}_C^n,\qquad
B^n(C)\coloneqq\operatorname{im}\mathrm{d}_C^{n-1}.
\]
记像分解为 \(\mathrm{d}_C^{n-1}=i^n\pi^{n-1}\)，其中 \(\pi^{n-1}:C^{n-1}\to B^n(C)\) 满，\(i^n:B^n(C)\to C^n\) 单。由 \(\mathrm{d}_C^ni^n\pi^{n-1}=0\) 及 \(\pi^{n-1}\) 满，得到 \(\mathrm{d}_C^ni^n=0\)。因此，核嵌入 \(z^n:Z^n(C)\to C^n\) 的泛性质唯一给出 \(j^n:B^n(C)\to Z^n(C)\)，使
\[
z^nj^n=i^n.
\]
因为 \(i^n\) 单，\(j^n\) 也单。这就是余边对象到余圈对象的典范嵌入，于是定义
\[
H^n(C)\coloneqq\operatorname{coker}\bigl(j^n:B^n(C)\longrightarrow Z^n(C)\bigr).
\]
它称为次数 \(n\in\mathbb Z\) 的\term[shangtongdiao]{上同调对象}[cohomology object]。当 \(R\) 为含幺环、\(C\) 为幺性左 \(R\) 模复形时，\(Z^n(C)\) 的元素称\term[yuquan]{余圈}[cocycle]，\(B^n(C)\) 的元素称\term[yubian]{余边}[coboundary]，且 \(H^n(C)=Z^n(C)/B^n(C)\)。对 \(\mathcal A\) 中的链复形 \((C_n,\partial_n)\)，相应定义 \(Z_n(C)=\ker\partial_n\)、\(B_n(C)=\operatorname{im}\partial_{n+1}\)，余核 \(H_n(C)\coloneqq\operatorname{coker}\bigl(B_n(C)\longrightarrow Z_n(C)\bigr)\) 称为次数 \(n\) 的\term[tongdiao]{同调对象}[homology object]；对模链复形，\(Z_n(C)\) 与 \(B_n(C)\) 的元素分别称\term[quan]{圈}[cycle]与\term[bianjie]{边界}[boundary]。
\end{definition}
\symidx{\(Z_n(C),B_n(C),H_n(C)\)：链复形的圈、边界与同调对象}
\symidx{\(Z^n(C),B^n(C),H^n(C)\)：余圈、余边与上同调对象}

在模范畴中，两个余圈代表同一上同调类，恰好表示它们之差是余边：能够由前一步微分产生的改变被忽略了。在一般交换范畴中，相同的思想由上述余核表达，不能预先假定对象带有元素。

\begin{definition}\label{b4:def:acyclic-complex}
设 \(C\) 是交换范畴中的复形。若所有次数的同调对象均为零，则称 \(C\) 为\term[lingtiao-fuxing]{零调复形}[acyclic complex]，也称正合复形，即在每一个次数处都正合。对上链复形，这等价于 \(\operatorname{im}\mathrm{d}_C^{n-1}=\ker\mathrm{d}_C^n\) 对每个 \(n\in\mathbb Z\) 成立，其中等号指同一个子对象；链复形作相应的指标转换。
\end{definition}

\begin{example}\label{b4:exm:integer-three-complex}
取交换群复形
\[
C^0=\mathbb Z\xrightarrow{\ \mathrm{d}^0\ }C^1=\mathbb Z^2
\xrightarrow{\ \mathrm{d}^1\ }C^2=\mathbb Z,
\qquad \mathrm{d}^0(a)=(2a,0),\quad \mathrm{d}^1(x,y)=3y.
\]
复合为零，且
\[
\begin{aligned}
H^0(C)&=0,\\
H^1(C)&=\mathbb Z(1,0)/2\mathbb Z(1,0)\cong\mathbb Z/2\mathbb Z,\\
H^2(C)&=\mathbb Z/3\mathbb Z.
\end{aligned}
\]
中间的核由 \((1,0)\) 生成，前一步只得到其偶数倍；末项没有下一步限制，却须除去 \(3\mathbb Z\)。这两个商分别记录不同次数的非正合性。
\end{example}

\subsection{指标转换与次数平移约定}
\label{b4:subsec:0501:03}

把链复形改记为上链复形时，本书取
\[
C^n=C_{-n},\qquad \mathrm{d}_C^n=\partial_{-n},\qquad H^n(C)=H_{-n}(C).
\]
因为 \(C^{n+1}=C_{-(n+1)}=C_{-n-1}\)，原来的 \(\partial_{-n}:C_{-n}\to C_{-n-1}\) 正好成为 \(\mathrm{d}_C^n:C^n\to C^{n+1}\)。这只是重编号，没有额外的负号；负号出现在下面的移位中。

\begin{definition}\label{b4:def:complex-shift}
设 \(C\) 是交换范畴中的上链复形，\(r\in\mathbb Z\)。规定
\[
C[r]^n=C^{n+r},\qquad \mathrm{d}_{C[r]}^n=(-1)^r\cdot\mathrm{d}_C^{n+r},
\]
所得复形称为 \(C\) 的 \(r\) 次\term[yiwei]{移位}[shift]。对复形之间的零次态射，移位只改指标，不增加符号。
\end{definition}
\symidx{\(C[r]\)：复形的移位，\(C[r]^n=C^{n+r}\)}

无论是否乘上这个整体符号，相邻微分的复合都仍为零；采用 \((-1)^r\) 并不是复形条件本身所强制的。核与像不因微分乘 \(-1\) 改变，所以
\[
H^n(C[r])=H^{n+r}(C),\qquad (C[r])[s]=C[r+s].
\]
原来次数为 \(m\) 的对象，在 \(C[r]\) 中出现于次数 \(m-r\)，因为 \(C[r]^{m-r}=C^m\)。特别地，\(C[1]^{-1}=C^0\)、\(C[1]^0=C^1\)，所以 \([1]\) 把原来的各项向较小次数移动一格。例如 \(M[0][1]\) 的唯一非零项在次数 \(-1\)，不是次数 \(1\)。若沿用链记号，同一约定写成 \((C[r])_n=C_{n-r}\)、\(\partial_{C[r]}=(-1)^r\partial_C\)。阅读采用相反移位记号的文献时，须同时转换次数和微分。

\ProseParagraphBreak
移位的负号使复形能够相容地组成映射锥，并使同伦在移位后仍满足同一形式的公式。此处先固定约定，\secref{b4:sec:0602}将通过微分平方和连接同态的计算说明它的作用。
